\section{Genetic algorithms: a concise historical arc}
\label{sec:ga}

Holland's framework emphasized string representations, crossover, and mutation with schema-intuitive explanations~\cite{holland1975adaptation}.
Goldberg systematized algorithm engineering practice~\cite{goldberg1989genetic}.
Parallel communities developed evolution strategies (ES) emphasizing continuous parameter spaces and self-adaptation~\cite{hansen2001completely,beyer2002evolution}.

\subsection{Modern contrasts}
\begin{itemize}
  \item \textbf{Multi-objective dominance}: NSGA-II remains a widely deployed baseline~\cite{deb2002fast}.
  \item \textbf{Distribution-based search}: CMA-ES adapts full covariance of search distributions~\cite{hansen2001completely}.
  \item \textbf{Quality diversity}: MAP-Elites archives elites across behavioral niches~\cite{mouret2015illuminating}.
\end{itemize}

\begin{figure}[t]
  \centering
  \begin{tikzpicture}[
    font=\footnotesize,
    agent/.style={circle, draw, minimum size=0.55cm},
    arr/.style={-{Stealth[length=1.8mm]}}
  ]
  \node[agent, fill=blue!15] (r) at (0,0) {$r$};
  \node[agent, fill=blue!10] (a1) at (-1.5,-1.2) {$a_1$};
  \node[agent, fill=blue!10] (a2) at (1.5,-1.2) {$a_2$};
  \node[agent, fill=orange!20] (b1) at (-2.2,-2.6) {$b_1$};
  \node[agent, fill=orange!15] (b2) at (-0.8,-2.6) {$b_2$};
  \draw[arr] (r) -- (a1);
  \draw[arr] (r) -- (a2);
  \draw[arr] (a1) -- (b1);
  \draw[arr] (a1) -- (b2);
  \node[align=left] at (3,-1.3) {Archive retains\\multiple lineages};
\end{tikzpicture}

  \caption{Toy illustration of an evolutionary archive as a tree of edits (original TikZ sketch).}
  \label{fig:archive-tree}
\end{figure}

\subsection{Transition to code evolution}
When genotypes become programs in Turing-complete languages, syntactic constraints and semantic brittleness dominate~\cite{koza1992genetic}.
LLM-based mutation operators bias proposals toward human-like code edits~\cite{lange2025shinka}.
